\subsection{EXAMPLE 5.2: INDEPENDENCE}

\textbf{Enunciado}

\emph{A fundamental assumption of the random walk hypothesis of market
returns is that returns from one time period to the next are statistically
independent. This assumption implies}
\begin{enumerate}[label=\alph*.]
\item Returns from one time period to the next can never be equal.
\item Returns from one time period to the next are uncorrelated.
\item Knowledge of the returns from one time period does not help in
predicting returns from the next time period.
\item \hl{Both b. and c. are true.}
\end{enumerate}

\subsubsection*{Solución}


\textbf{Definición de independencia. }Sean \(R_t\) y \(R_{t+1}\) rendimientos consecutivos. Son \textbf{independientes} si:
\[
\Prob(R_t\in A,\;R_{t+1}\in B)=\Prob(R_t\in A)\Prob(R_{t+1}\in B)
\quad \forall A,B.
\]
Intuitivamente: conocer \(R_t\) no cambia la distribución de \(R_{t+1}\).

\textbf{Implicación sobre correlación (b). }Si \(R_t\) y \(R_{t+1}\) son independientes y con varianza finita, entonces:
\[
\mathrm{Cov}(R_t,R_{t+1})
=\E[R_tR_{t+1}]-\E[R_t]\E[R_{t+1}]
=0,
\]
pues por independencia \(\E[R_tR_{t+1}]=\E[R_t]\E[R_{t+1}]\).
(b) es verdadera.

\textbf{Implicación predictiva (c). }Independencia implica:
\[
\Prob(R_{t+1}\in B\mid R_t\in A)=\Prob(R_{t+1}\in B),
\]
y por tanto \(R_t\) no aporta información para predecir \(R_{t+1}\).
(c) es verdadera.

\textbf{Por qué (a) es falsa.} Independencia \emph{no} prohíbe que \(R_{t+1}=R_t\). En distribuciones discretas puede ocurrir con probabilidad positiva, y en continuas puede ocurrir con probabilidad 0 por continuidad (no por independencia).
(a) es falsa.


Como (b) y (c) son verdaderas:
(d) Both b. and c. are true.












